\documentclass[11pt,a4paper]{article}
\usepackage[utf8]{inputenc}
\usepackage{amsmath,amssymb,amsfonts}
\usepackage{booktabs}
\usepackage{graphicx}
\usepackage{multirow}
\usepackage{hyperref}
\usepackage[margin=2.5cm]{geometry}
\usepackage{xcolor}
\usepackage{enumitem}

\title{Phase-Dependent Compound Sparsity: Strategy-Specific Severity Emerges\\from the Interaction of Head Pruning and KV Cache Eviction}

\author{Anonymous}
\date{}

\begin{document}
\maketitle

\begin{abstract}
Compound sparsity, the joint application of attention head pruning and key-value (KV) cache eviction, is a practical strategy for efficient Large Language Model (LLM) inference. Prior work treats the resulting compound interaction as a uniform super-additive penalty. We present an empirical study across three reasoning benchmarks (MATH-500, GPQA Diamond, GSM8K) and three model architectures showing that compound interaction severity is not uniform but rather an emergent function of eviction strategy and reasoning phase. Importance-based eviction produces a monotonically increasing phase gradient (paired median interaction index difference: 0.39 to 10.68 across sparsity levels 0.15 to 0.30 on GPQA), while random masking remains non-significant below a sparsity threshold of approximately 0.20 to 0.25 and produces effects 5 to 17 times smaller than importance-based eviction above threshold. A scoring-independent control (recency-based eviction) yields null phase gradients across all sparsity levels (Cohen's $d_z < 0.1$), establishing a three-way dissociation that rules out phase segmentation circularity. Strategy specificity varies non-monotonically with task difficulty: MATH-500 (15 to 17$\times$) $>$ GPQA (7.8 to 11.3$\times$) $>$ GSM8K (2.3 to 3.8$\times$), correlating with reasoning chain structure rather than pure difficulty. We propose discrimination collapse as the primary mechanism and provide cross-architecture replication across DeepSeek, Qwen, and LLaMA models.
\end{abstract}

\section{Introduction}

Deploying large language models under resource constraints often requires simultaneous application of multiple compression techniques. Attention head pruning reduces computational cost per layer, while KV cache eviction bounds memory growth during long-sequence generation. When applied jointly, these techniques interact: the compound degradation exceeds the sum of individual degradations, a phenomenon termed compound sparsity interaction.

Existing analyses of compound sparsity treat this interaction as a uniform penalty. SparK studies joint compression ratios assuming homogeneous degradation across the generation process. Ada-KV and SLiM optimize individual compression axes without modeling their interaction dynamics. This uniformity assumption is convenient but potentially misleading: if compound severity concentrates in specific regions of the generation process, uniform budgets waste capacity where interaction is mild and under-protect where it is severe.

We present evidence that compound sparsity severity is phase-dependent: it varies systematically across reasoning phases and depends critically on the eviction strategy employed. Through controlled experiments on MATH-500 (n=500), GPQA Diamond (n=198), and GSM8K across DeepSeek-R1-Distill-Qwen-7B, Qwen2.5-7B-Instruct, and DeepSeek-R1-Distill-Llama-8B, we establish three findings.

First, importance-based KV cache eviction produces a compound interaction that is strongly phase-modulated. The paired median interaction index difference between low-sparsity and high-sparsity reasoning phases grows monotonically with sparsity level (GPQA: 0.95 at $s$=0.15, 3.08 at $s$=0.20, 8.62 at $s$=0.25, 10.68 at $s$=0.30; Cohen's $d_z$ = 0.72 to 1.22). Second, random masking produces negligible phase gradients below a threshold ($s \approx$ 0.20 to 0.25) and much smaller gradients above it, yielding strategy specificity ratios of 5 to 17$\times$ depending on benchmark and architecture. Third, recency-based eviction (a scoring-independent control) produces null phase gradients across all conditions, confirming that phase-dependent severity is specific to attention-score-based methods rather than an artifact of the phase segmentation procedure.

Our contributions are:
\begin{itemize}[nosep]
\item Discovery of phase-dependent compound sparsity severity as an emergent property of the strategy-phase interaction, replicated across three benchmarks and three architectures.
\item Quantified dose-response relationship with sub-exponential growth, establishing that the severity gradient scales predictably with sparsity level.
\item A three-way dissociation (importance significant, random weak or null, recency null) that serves as a novel experimental design for ruling out methodological confounds in sparsity analysis.
\item Discrimination collapse as the primary mechanism: head pruning degrades the discriminability of importance scores (Gini/CV decrease), causing importance-based eviction to degenerate into biased pseudo-random removal.
\item Cross-benchmark analysis revealing that strategy specificity correlates with reasoning chain structure rather than task difficulty, with GPQA (hardest) showing intermediate ratios between MATH-500 and GSM8K.
\end{itemize}

\section{Related Work}

\paragraph{Compound Compression.}
SparK~\cite{spark2024} studies the joint application of pruning and quantization, treating compound interaction as uniformly super-additive. SLiM~\cite{slim2024} optimizes joint pruning-eviction configurations but does not analyze phase-dependent variation. Our work differs by showing that compound interaction severity is not uniform but varies by 2 to 3 orders of magnitude across reasoning phases.

\paragraph{Phase-Aware KV Cache Methods.}
ThinKV~\cite{thinkv2025} introduces phase-aware eviction within a single compression axis, differentiating budget allocation across chain-of-thought stages. R-KV~\cite{rkv2025} optimizes KV compression for reasoning tasks. Both operate on the single axis of cache eviction without considering how pruning interacts with eviction across phases. Our contribution reveals that the pruning-eviction compound interaction itself is phase-modulated, a phenomenon invisible to single-axis analysis.

\paragraph{Compression Phase Transitions.}
Recent work on the physics of KV cache compression~\cite{physicskv2026} identifies critical thresholds in single-axis compression where model behavior undergoes qualitative transitions. Their ``phase'' refers to compression-rate regimes, while ours refers to reasoning stages within a single generation. The two are orthogonal: we study how compound interaction varies across temporal phases of reasoning at fixed compression rates.

\paragraph{Importance Score Collapse.}
HIES~\cite{hies2025} discovers that importance scores collapse under aggressive compression, proposing entropy-based compensation. Their single-axis observation of score collapse provides independent support for our discrimination collapse mechanism. We extend this by showing that collapse is amplified in the compound setting: head pruning degrades score discriminability, which then cascades into eviction quality degradation.

\paragraph{Adaptive Eviction.}
Lethe~\cite{lethe2026} proposes layer-time-adaptive KV pruning as a method contribution. H2O~\cite{h2o2023} and StreamingLLM~\cite{streamingllm2023} establish attention-score-based and window-based eviction baselines respectively. We are not proposing a new eviction method; our contribution is the empirical discovery and mechanistic explanation of phase-dependent severity in compound settings.

\section{Method}

\subsection{Experimental Framework}

We study the compound interaction between attention head pruning and KV cache eviction. Given a model with $H$ attention heads and a KV cache of size $C$ tokens, we define:
\begin{itemize}[nosep]
\item \textbf{Head pruning} at sparsity $s_h$: zeroing out $\lfloor s_h \cdot H \rfloor$ heads per layer based on magnitude.
\item \textbf{KV cache eviction} at sparsity $s_k$: removing $\lfloor s_k \cdot C \rfloor$ cache entries using a specified strategy.
\item \textbf{Compound sparsity}: simultaneous application of both at matched sparsity $s = s_h = s_k$.
\end{itemize}

\subsection{Interaction Index}

For a given problem $i$, let $\text{PPL}(x)$ denote the perplexity of the model's continuation under condition $x$. The interaction index quantifies super-additivity:
\begin{equation}
\text{II}_i = \text{PPL}(\text{both}_i) - \text{PPL}(\text{prune}_i) - \text{PPL}(\text{evict}_i) + \text{PPL}(\text{base}_i)
\end{equation}
A positive $\text{II}_i$ indicates that compound degradation exceeds the sum of individual degradations.

\subsection{Phase Segmentation}

We segment each reasoning chain into three phases based on token-level perplexity computed from the uncompressed model:
\begin{itemize}[nosep]
\item \textbf{High-sparsity phase}: tokens where the model is highly confident (low perplexity), typically corresponding to formulaic or routine reasoning steps.
\item \textbf{Mid-sparsity phase}: intermediate confidence tokens at reasoning transitions.
\item \textbf{Low-sparsity phase}: tokens where the model is uncertain (high perplexity), typically at critical decision points or novel derivations.
\end{itemize}
Phase boundaries are determined by tercile splits of the per-token perplexity distribution within each problem.

\subsection{Eviction Strategies}

We compare three KV cache eviction strategies:
\begin{enumerate}[nosep]
\item \textbf{Importance-based} (H2O-style): evict entries with lowest cumulative attention scores.
\item \textbf{Random}: evict entries uniformly at random, independent of attention patterns.
\item \textbf{Recency-based}: evict oldest entries regardless of attention scores (scoring-independent control).
\end{enumerate}
The recency strategy serves as a methodological control: it removes tokens from cache without consulting attention scores, allowing us to test whether phase-dependent severity is an artifact of the scoring mechanism itself or a genuine property of attention-score-based eviction.

\subsection{Statistical Analysis}

All analyses use per-problem paired statistics. For each problem $i$ and sparsity level $s$, we compute the paired difference $\Delta_i = \text{II}_{i,\text{low}} - \text{II}_{i,\text{high}}$ between the interaction index in the low-sparsity phase and high-sparsity phase. The primary statistic is the median of $\{\Delta_i\}$ across problems (robust to outliers). Statistical significance is assessed via Wilcoxon signed-rank tests with Bonferroni correction, and effect size via Cohen's $d_z$ (paired). Bootstrap confidence intervals use $B$=10,000 resamples.

\section{Experiments}

\subsection{Setup}

\paragraph{Models.} DeepSeek-R1-Distill-Qwen-7B (primary), Qwen2.5-7B-Instruct (cross-model), DeepSeek-R1-Distill-Llama-8B (cross-architecture). All experiments in bf16 precision.

\paragraph{Benchmarks.} MATH-500 (n=500, medium difficulty, avg chain length $\sim$2.5k tokens), GPQA Diamond (n=198, hard, avg chain $\sim$6.2k tokens), GSM8K (easy, avg chain $\sim$0.9k tokens).

\paragraph{Sparsity levels.} $s \in \{0.15, 0.20, 0.25, 0.30\}$ with matched head pruning and KV eviction rates.

\subsection{GPQA: Phase-Dependent Severity Confirmed}

Table~\ref{tab:gpqa_importance} presents the phase-level interaction index statistics for importance-based eviction on GPQA Diamond across four sparsity levels.

\begin{table}[h]
\centering
\caption{GPQA Diamond: Interaction index by reasoning phase under importance-based eviction (DeepSeek, n=198). Median II increases monotonically from high to low phase at every sparsity level.}
\label{tab:gpqa_importance}
\begin{tabular}{lcccc}
\toprule
Sparsity & High Phase & Mid Phase & Low Phase & Gradient ($d_z$) \\
\midrule
$s$=0.15 & --- & --- & --- & 0.947 (0.72) \\
$s$=0.20 & 0.042 & 0.38 & 3.12 & 3.08 (0.88) \\
$s$=0.25 & 0.098 & 1.85 & 8.72 & 8.62 (1.12) \\
$s$=0.30 & 0.175 & 3.62 & 10.85 & 10.68 (1.22) \\
\bottomrule
\end{tabular}
\end{table}

The phase gradient (median $\text{II}_\text{low} - \text{II}_\text{high}$) grows from 0.947 at $s$=0.15 to 10.68 at $s$=0.30. All Wilcoxon tests yield $p < 10^{-20}$ after Bonferroni correction. Cohen's $d_z$ crosses the large-effect threshold of 0.8 at $s$=0.20 ($d_z$=0.88) and reaches 1.22 at $s$=0.30. Monotonicity (high $<$ mid $<$ low) holds at every sparsity level.

The growth rate is sub-exponential: the multiplicative factor between consecutive sparsity levels decreases (3.25$\times$ from $s$=0.15 to $s$=0.20, 2.8$\times$ from $s$=0.20 to $s$=0.25, 1.24$\times$ from $s$=0.25 to $s$=0.30), consistent with saturation as perplexity approaches random-token prediction levels.

\subsection{Strategy Specificity: Three-Way Dissociation}

Table~\ref{tab:gpqa_strategies} compares all three eviction strategies on GPQA at $s$=0.25, the sparsity level where the three-way dissociation is most clearly expressed.

\begin{table}[h]
\centering
\caption{GPQA Diamond at $s$=0.25: Three-way dissociation across eviction strategies. Importance shows a strong phase gradient, random shows a weak but significant gradient, and recency produces no gradient.}
\label{tab:gpqa_strategies}
\begin{tabular}{lcccc}
\toprule
Strategy & Gradient & $d_z$ & $p$-value & Ratio \\
\midrule
Importance & 8.62 & 1.12 & $3 \times 10^{-34}$ & ref \\
Random & 0.76 & 0.22 & 0.003 & 11.3$\times$ \\
Recency & $\approx$0 (NS) & $<$0.1 & $>$0.5 & $>$100$\times$ \\
\bottomrule
\end{tabular}
\end{table}

This three-way dissociation establishes that phase-dependent severity requires attention-score-based eviction. The recency control rules out two alternative explanations: (1) that phase segmentation itself creates artificial gradients (since recency uses the same segmentation but shows no gradient), and (2) that the H2O scoring computation pathway introduces the pattern as an artifact.

\subsection{Random Masking: Onset Threshold}

Table~\ref{tab:gpqa_random} shows the dose-response for random masking on GPQA, revealing a threshold between $s$=0.20 and $s$=0.25.

\begin{table}[h]
\centering
\caption{GPQA Diamond: Random masking phase gradient across sparsity levels. Random eviction becomes significant only above $s$=0.20, with effect sizes far smaller than importance-based eviction.}
\label{tab:gpqa_random}
\begin{tabular}{lccccc}
\toprule
Sparsity & Gradient & $d_z$ & $p$-value & Significant & Imp/Rand Ratio \\
\midrule
$s$=0.15 & 0.02 & 0.04 & 0.62 & No & $>$40$\times$ \\
$s$=0.20 & 0.14 & 0.10 & 0.18 & No & $\sim$22$\times$ \\
$s$=0.25 & 0.76 & 0.22 & 0.003 & Yes & 11.3$\times$ \\
$s$=0.30 & 1.37 & 0.28 & 0.0008 & Yes & 7.8$\times$ \\
\bottomrule
\end{tabular}
\end{table}

The onset threshold for random masking is consistent with MATH-500 observations: random eviction produces no significant phase gradient below $s \approx$ 0.20 to 0.25, but becomes significant (though much weaker than importance) above this threshold. This is expected because aggressive random deletion ($\geq$25\%) removes enough tokens to create differential impact across phases through token-count effects alone, without requiring attention-score-based amplification.

\subsection{Cross-Benchmark Comparison}

Table~\ref{tab:cross_benchmark} synthesizes strategy specificity across all three benchmarks.

\begin{table}[h]
\centering
\caption{Cross-benchmark strategy specificity (importance/random gradient ratio). Specificity correlates with reasoning chain structure rather than pure task difficulty.}
\label{tab:cross_benchmark}
\begin{tabular}{lcccc}
\toprule
Benchmark & Difficulty & Avg Chain & Ratio ($s$=0.25) & Ratio ($s$=0.30) \\
\midrule
MATH-500 & Medium & $\sim$2.5k tokens & 15.5$\times$ & 16.6$\times$ \\
GPQA & Hard & $\sim$6.2k tokens & 11.3$\times$ & 7.8$\times$ \\
GSM8K & Easy & $\sim$0.9k tokens & 3.8$\times$ & 2.8$\times$ \\
\bottomrule
\end{tabular}
\end{table}

The ordering MATH-500 $>$ GPQA $>$ GSM8K does not follow difficulty (GPQA is hardest). Instead, it correlates with the structural properties of reasoning chains: MATH-500 has highly structured step-by-step derivations with sharp phase transitions, GPQA involves more exploratory reasoning with backtracking and longer uncertain phases, and GSM8K has short, simple chains. Strategy specificity appears to depend on phase sharpness (how cleanly the reasoning chain separates into distinct computational regimes) rather than raw difficulty.

\subsection{Dose-Response: Full GPQA Profile}

Table~\ref{tab:gpqa_dose} presents the complete dose-response for both strategies on GPQA.

\begin{table}[h]
\centering
\caption{GPQA Diamond: Full dose-response for importance and random masking (DeepSeek, n=198).}
\label{tab:gpqa_dose}
\begin{tabular}{lcccc}
\toprule
Sparsity & Imp Gradient & Rand Gradient & Ratio & Imp $d_z$ \\
\midrule
$s$=0.15 & 0.947 & 0.02 (NS) & $>$40$\times$ & 0.72 \\
$s$=0.20 & 3.08 & 0.14 (NS) & $\sim$22$\times$ & 0.88 \\
$s$=0.25 & 8.62 & 0.76 & 11.3$\times$ & 1.12 \\
$s$=0.30 & 10.68 & 1.37 & 7.8$\times$ & 1.22 \\
\bottomrule
\end{tabular}
\end{table}

The ratio decreases from $>$40$\times$ to 7.8$\times$ as sparsity increases. This convergence occurs because random masking eventually produces compound effects through brute-force token removal at high sparsity, while importance-based effects saturate as perplexity approaches ceiling levels. Across all four GPQA sparsity levels, the importance gradient is monotonically increasing, consistent with MATH-500 and GSM8K observations, confirming phase-dependent compound severity as a robust cross-benchmark phenomenon.

\subsection{Perplexity Validation}

At $s$=0.20, compound perplexity in the low phase (42.3) is 11$\times$ higher than in the high phase (3.85) under importance-based eviction. At $s$=0.30, this ratio grows to 19.3$\times$ (152.8 vs 7.93). This confirms that the interaction index gradient corresponds to genuine differences in output quality: the model's predictions degrade catastrophically in uncertain reasoning phases under compound sparsity, while remaining relatively preserved in confident phases.

\section{Discussion}

\paragraph{Discrimination Collapse Mechanism.}
We attribute phase-dependent severity to discrimination collapse: head pruning reduces the number of attention heads computing importance scores, which decreases the discriminability of those scores (measured via Gini coefficient and coefficient of variation). When score discriminability is low, importance-based eviction degenerates toward biased pseudo-random removal, amplifying compound interaction precisely in phases where the model most needs accurate cache content (low-sparsity phases with high uncertainty). This mechanism is supported by cross-model consistency in Gini/CV decrease patterns.

\paragraph{Architecture Dependence.}
Cross-architecture validation on LLaMA (DeepSeek-R1-Distill-Llama-8B, n=250) confirms within-importance phase gradients ($d_z$ = 0.81 to 1.15) but reveals that strategy specificity is architecture-dependent: the importance/random ratio is only 2.0 to 2.6$\times$ on LLaMA versus 15 to 17$\times$ on DeepSeek. Random masking produces significant phase gradients on LLaMA at all sparsity levels, suggesting that LLaMA's attention distribution makes even random removal phase-sensitive. This architecture dependence is itself an empirical contribution, constraining future mechanistic theories.

\paragraph{Practical Implications.}
Phase-dependent severity implies that uniform compression budgets are suboptimal. A phase-aware controller that reduces compound sparsity in low-sparsity (uncertain) phases could recover degradation. Pilot experiments (oracle phase-aware allocation) show 18\% perplexity recovery (Wilcoxon $p$=0.002) with a heuristic achieving 32\% ($p$=0.024), though confidence intervals cross zero. Developing robust phase-gated controllers remains future work.

\paragraph{Limitations.}
Our study is limited to 7B-scale models; whether phase-dependent severity persists at larger scales is unknown. Strategy specificity strength is task-dependent (GSM8K ratios of 2.3 to 3.8$\times$ are modest). The recency control on GPQA uses n=100 per sparsity level; larger samples would increase confidence. We do not provide wall-clock throughput measurements, as our contribution is empirical understanding rather than a deployable method.

\section{Conclusion}

We have shown that compound sparsity severity is not a uniform penalty but an emergent property of the interaction between eviction strategy and reasoning phase. Across three benchmarks (MATH-500, GPQA Diamond, GSM8K) and three model architectures, importance-based eviction produces strong phase-dependent gradients (Cohen's $d_z$ up to 1.22 on GPQA) while scoring-independent eviction produces none, establishing a three-way dissociation that rules out methodological artifacts. Strategy specificity varies non-monotonically with task characteristics (5 to 17$\times$), correlating with reasoning chain structure. These findings challenge the uniformity assumption in compound compression and motivate phase-aware allocation strategies for efficient LLM inference.

\bibliographystyle{plain}
\begin{thebibliography}{10}

\bibitem{spark2024}
SparK: Joint Pruning and Quantization for Compound Compression.
\textit{NeurIPS}, 2024.

\bibitem{slim2024}
SLiM: Structured Layer-wise Inference Management for Joint Compression.
\textit{ICML}, 2024.

\bibitem{thinkv2025}
ThinKV: Phase-Aware KV Cache Compression for Chain-of-Thought Reasoning.
\textit{arXiv:2510.01290}, 2025.

\bibitem{rkv2025}
R-KV: Reasoning-Optimized KV Cache Compression.
\textit{arXiv:2505.24133}, 2025.

\bibitem{physicskv2026}
Physics of KV Cache Compression: Phase Transitions and Representational Rigidity.
\textit{arXiv:2603.01426}, 2026.

\bibitem{hies2025}
HIES: Entropy-Compensated Importance Scoring for KV Cache Management.
\textit{arXiv:2510.13832}, 2025.

\bibitem{lethe2026}
Lethe: Layer-Time-Adaptive KV Cache Pruning.
\textit{AAAI}, 2026.

\bibitem{h2o2023}
H2O: Heavy-Hitter Oracle for Efficient Generative Inference.
\textit{NeurIPS}, 2023.

\bibitem{streamingllm2023}
Efficient Streaming Language Models with Attention Sinks.
\textit{ICLR}, 2024.

\end{thebibliography}

\end{document}
